Playtesting is essential for understanding how game rules influence player choices and outcomes. Human testing provides valuable feedback, but repeatedly evaluating possible rule changes requires time and participants. Automated agents can complement that process by generating repeatable games and quantitative observations. Previous work has used agents with different styles to study board-game configurations and has used self-play to investigate alternative chess rules \cite{silva2017,tomasev2020}.

The research question for this project is: How do small Backgammon rule changes affect the relative performance and observable behavior of Monte Carlo agents with different heuristic biases? The aim is to build an experimental environment that can expose changes in strategic incentives, rather than to produce an expert-strength Backgammon player.

Four intended styles are rear-priority, front-priority, aggressive, and defensive. They will respectively favor advancing trailing checkers, advancing leading checkers, hitting opposing blots, and protecting checkers or building safe positions. All four will share the same underlying decision method and configured search budget. Before introducing rule variations, their intended behavioral differences must be measured under standard rules.

At midterm, the main achievements are an operational Backgammon engine, automated rule and state checks, configurable flat Monte Carlo search, a shared cutoff heuristic, and experiments comparing random and heuristic-guided rollout policies. These establish the platform and initial baseline needed for the central research experiments. No conclusions about rule variations or the four planned styles have been obtained yet.
